%\documentclass{beamer}
%\RequirePackage[2020-02-02]{latexrelease}
\documentclass[hyperref={pdfpagemode=FullScreen},aspectratio=169]{beamer}

%    1610: 16:10
%    169: 16:9
%    149: 14:9
%    141: 1.41:1
%    54: 5:4
%    43: 4:3 [default]
%    32: 3:2

\setbeamersize{text margin left=1pt,text margin right=1pt}
\mode<presentation> {
%\setbeamertemplate{caption}[numbered]
\setbeamertemplate{footline}[page number]
\setbeamertemplate{bibliography item}{\insertbiblabel}
\usetheme{default}
\usefonttheme{serif}
}
\usepackage{amsfonts,epsfig,subcaption,graphicx,tabularx,amsmath,bm,relsize,steinmetz,booktabs,epstopdf}
\usepackage{rotating}
\usepackage{extarrows,mathtools}
\usepackage{amsrefs}
\usepackage[english]{babel}
\usepackage{epstopdf}
\usepackage{listings}
\usepackage{color}
\usepackage{mathrsfs}
\usepackage{tcolorbox}
\usepackage{arydshln}
\setlength{\dashlinedash}{4pt}
\setlength{\dashlinegap}{4pt}

\usepackage[nameinlink]{cleveref}
\usepackage{aliascnt}
\crefname{equation}{Eqn.}{Eqns.}
\newaliascnt{ineq}{equation}
\aliascntresetthe{ineq}
\crefname{ineq}{Ineq.}{Ineqs.}
\creflabelformat{ineq}{#2{\upshape(#1)}#3} 
\makeatletter
\def\ineq{$$\refstepcounter{ineq}}
\def\endineq{\eqno \hbox{\@eqnnum}$$\@ignoretrue}
\makeatother
\crefname{figure}{Fig.}{Figs.}
\crefname{table}{Table}{Tables}
%\renewcommand{\figurename}{Fig.} %without babel package
\addto\captionsenglish{\renewcommand{\figurename}{Fig.}} %with babel package
%
%\usepackage{amsthm}
%\newtheorem{definition}{Definition}

\usepackage[printwatermark]{xwatermark}
\usepackage{xcolor}
\usepackage{graphicx}
\usepackage{tikz}
\usetikzlibrary{patterns,decorations.pathmorphing,decorations.markings}
\usetikzlibrary{shapes,arrows,positioning,calc}
\usetikzlibrary{patterns,angles,quotes}
\usepackage{circuitikz}

\newsavebox\mybox
\savebox\mybox{\tikz[color=red,opacity=0.1]\node{\small{Dr. Y. V. Karteek, SCEE, IIT Mandi}};}
\newwatermark*[
  allpages,
  angle=0,
  scale=1,
  xpos=-50,
  ypos=-38
]{\usebox\mybox}

%\title{\Large{UEI303: Techniques on Signals and Systems \\~\\ Random Signals}~\\}
%\title{\Large{UEI407: Signals and Systems\\~\\ Random Signals}~\\}
\title{\Large{EE538: Data-driven Control\\~\\}~\\}
\author{\small Dr. Y. ~V. ~Karteek}
 
\date{} 
%\date{\tiny {\today}}
%----------------------------------------------------
\begin{document}
%----------------------------------------------------
\begin{frame}[plain]
          \begin{center}
                    \Large{}
                    \end{center}
               \titlepage
        \end{frame}
%--------------------------------------------------


\begin{frame}{Motivation}
\begin{itemize}
\item Many real systems lack accurate mathematical models
\item Nonlinear systems are difficult to model and control
\item Traditional approach:
\begin{itemize}
\item Identify model $\rightarrow$ design controller
\end{itemize}
\item VRFT approach:
\begin{itemize}
\item Directly design controller from data
\item No plant model required
\end{itemize}
\end{itemize}
\end{frame}

%------------------------------------------------

\begin{frame}{Key Idea of VRFT}
\begin{itemize}
\item Use input-output data from the plant
\item Define a \textbf{virtual reference}
\item Convert control problem into a data-driven optimization
\item One-shot design (no iterations required)
\end{itemize}
\end{frame}

%------------------------------------------------

\begin{frame}{Control System Setup}
\begin{itemize}
\item Nonlinear plant:
\begin{align}
y(t) = P(u(t))
\end{align}
\item Controller:
\begin{align}
u(t) = C_\theta(e(t)), \quad e(t) = r(t) - y(t)
\end{align}
\item Goal: match reference model $M$
\end{itemize}
\end{frame}

%------------------------------------------------

\begin{frame}{Control Objective}
\begin{itemize}
\item Desired behavior:
\begin{align}
y \approx M(r)
\end{align}
\item Optimization problem:
\begin{align}
\min_\theta | y_\theta - M(r) |^2
\end{align}
\item Problem: depends on unknown plant $P$
\end{itemize}
\end{frame}

%------------------------------------------------

\begin{frame}{Virtual Reference Concept}
\begin{itemize}
\item Construct virtual reference:
\begin{align}
\tilde{r} = M^{-1}(y)
\end{align}
\item Virtual error:
\begin{align}
\tilde{e} = \tilde{r} - y
\end{align}
\item Key idea:
\begin{itemize}
\item Use measured data $(u, y)$
\item Avoid explicit plant model
\end{itemize}
\end{itemize}
\end{frame}

%------------------------------------------------

\begin{frame}{VRFT Cost Function}
\begin{itemize}
\item Define data-based cost:
\begin{align}
J_{VRFT}(\theta) = | u - C_\theta(\tilde{e}) |^2
\end{align}
\item Depends only on data
\item No plant model needed
\end{itemize}
\end{frame}

%------------------------------------------------

\begin{frame}{Key Result}
\begin{block}{Theorem}
If perfect matching is achievable, minimizing $J_{VRFT}$ yields the same controller as the original control objective.
\end{block}

\begin{itemize}
\item VRFT replaces model-based optimization
\item Works directly with data
\end{itemize}
\end{frame}

%------------------------------------------------

\begin{frame}{Role of the Filter}
\begin{itemize}
\item When perfect matching is not possible:
\begin{itemize}
\item Introduce filter $F$
\end{itemize}
\item Modified cost:
\begin{align}
J_{VRFT}(\theta) = | F(u - C_\theta(\tilde{e})) |^2
\end{align}
\item Ensures:
\begin{itemize}
\item Approximate equivalence with true objective
\end{itemize}
\end{itemize}
\end{frame}

%------------------------------------------------

\begin{frame}{Filter Design}
\begin{itemize}
\item Optimal filter:
\begin{align}
F = (I - M) \frac{\partial P}{\partial u}
\end{align}
\item Interpretation:
\begin{itemize}
\item Matches curvature of cost functions
\item Improves performance in practical cases
\end{itemize}
\end{itemize}
\end{frame}

%------------------------------------------------

\begin{frame}{Advantages of VRFT}
\begin{itemize}
\item One-shot design (no iterations)
\item No plant model required
\item Works for nonlinear systems
\item Easy implementation for linear-in-parameter controllers
\end{itemize}
\end{frame}

%------------------------------------------------

\begin{frame}{Comparison with IFT}
\begin{itemize}
\item VRFT:
\begin{itemize}
\item One-shot
\item No gradient computation
\item No local minima issues
\end{itemize}
\item IFT (Iterative Feedback Tuning):
\begin{itemize}
\item Requires multiple experiments
\item Uses gradient descent
\end{itemize}
\end{itemize}
\end{frame}

%------------------------------------------------

\begin{frame}{Handling Noise}
\begin{itemize}
\item Real data contains noise
\item Solution:
\begin{itemize}
\item Use instrumental variables (IV)
\item Perform multiple experiments
\end{itemize}
\item Improves robustness of parameter estimation
\end{itemize}
\end{frame}

%------------------------------------------------

\begin{frame}{Multipass VRFT}
\begin{itemize}
\item If initial design is not sufficient:
\begin{itemize}
\item Apply controller
\item Collect new data
\item Reapply VRFT
\end{itemize}
\item Converges to desired behavior
\end{itemize}
\end{frame}

%------------------------------------------------

\begin{frame}{Summary}
\begin{itemize}
\item VRFT is a direct data-driven control design method
\item Uses virtual reference to bypass plant modeling
\item Effective for nonlinear systems
\item Practical and computationally efficient
\end{itemize}
\end{frame}

%------------------------------------------------

\begin{frame}{Thank You}
Questions?
\end{frame}

\end{document}
